\documentclass[10pt,journal]{IEEEtran}
\usepackage[T1]{fontenc}
\usepackage{amsmath,amssymb,booktabs,graphicx,cite,url}
\usepackage[hidelinks]{hyperref}
% Auto-generated scientific numbers for DLG-StreamMC R3
\newcommand{\PrimaryFOne}{0.610}
\newcommand{\LocalFOne}{0.600}
\newcommand{\DeltaFOne}{+0.010}
\newcommand{\PositiveSeeds}{3}
\newcommand{\NegativeSeeds}{2}
\newcommand{\StreamEvents}{100,000}
\newcommand{\UniqueContracts}{3,378}
\newcommand{\EventPrevalence}{3.10}
\newcommand{\EarlyTestContracts}{662}
\newcommand{\PrimaryLatency}{6.17}
\newcommand{\FullLatency}{8.07}
\newcommand{\LatencyReduction}{23.6}
\newcommand{\LossCount}{0}
\newcommand{\RestartErrors}{0}

\title{DLG-StreamMC: Bounded Stateful Local-to-Global Graph Inference with Validation-Calibrated Selective Execution}
\author{SeongSu Park and Ki-Hyung Kim}
\begin{document}
\maketitle
\begin{abstract}
Selective execution can reduce graph inference work, but a routing fraction is not a runtime measurement, and an event-weighted replay score is not contract-level detection accuracy. We study a bounded local-to-global system with a Graph Isomorphism Network fast path and selectively executed relational graph attention. One validation-frozen calibration, routing and fusion contract is shared by offline and event-driven inference. The primary policy obtains mean held-out contract F1 \PrimaryFOne{} versus \LocalFOne{} for the local path across five frozen models; the observed difference is descriptive. On a matched warmed event prefix, measured mean latency changes by \LatencyReduction{} percent relative to full relational execution. Both policies complete a retrospective replay of \StreamEvents{} events, with instrumented state and durable checkpoint restoration. An evidence audit distinguishes contract and event populations and identifies a global training-time boundary violation in the inherited evaluation artifacts. Consequently, the results support measured selective execution and bounded configured state on this workload, not strict live temporal detection, a confirmatory accuracy improvement, or constant process memory at arbitrary horizons.
\end{abstract}
\begin{IEEEkeywords}graph inference, selective execution, calibration, blockchain, stateful replay, reproducible evaluation\end{IEEEkeywords}

\section{Introduction}
Blockchain monitoring combines structural classification with a systems problem: local transaction neighborhoods evolve, model state occupies memory, and inference must produce a decision after each arriving event. A hierarchical detector can separate a relatively inexpensive local encoder from a relational stage. This separation offers a compute-allocation opportunity, but it also creates a reproducibility obligation. The score transformation, neighborhood construction, escalation rule and final threshold measured offline must be the ones executed in the profiler.

DLG-StreamMC extends the local-to-global decomposition of DLG-GNN~\cite{park2026decoupled} with bounded state and selective execution. The upstream multi-chain Graphs-of-Graphs representation~\cite{luo2024multichain} motivates contract-local graphs and inter-contract reasoning. The present study is a systems and evidence-integrity evaluation, not a new graph architecture. Its contribution is a shared execution interface, a measured accuracy--risk--cost description, and explicit limits on which empirical claims the available artifacts support.

The distinction matters in this revision. Historical offline, short-prefix and long-replay tables did not share identical score or relation semantics. We corrected the interface and reran inference rather than carrying their headline numbers forward. A further timestamp audit found that the retained cross-chain training pool was not globally earlier than every test contract, and that the entire replay prefix predates the latest training-reference observation. We expose this limitation instead of labeling retrospective execution as a leakage-free deployment experiment.

Our contributions are therefore: (i) policy-identical local and relational inference across separated evidence lanes; (ii) measured, bounded-state replay with explicit memory and restart diagnostics; (iii) validation-derived operating-point and dense risk sweeps; and (iv) prediction-paired descriptive statistics and a content-addressed evidence package. We do not claim universal superiority over graph, tabular or temporal models.

\section{Related Work and Positioning}
\subsection{Hierarchical and Temporal Graph Learning}
GIN studies graph-level representational power~\cite{xu2019gin}; GATv2 provides dynamic attention~\cite{brody2022how}. These components motivate the local and relational encoders, but do not establish the effectiveness of the selective system. The predecessor DLG-GNN decouples their execution. The predecessor remains an author-supplied manuscript; no unverified public identifier is assigned.

TGN~\cite{rossi2020temporal} and temporal attention~\cite{xu2020tgat} learn representations from timestamped histories. This is different from executing frozen encoders on incrementally materialized graph snapshots. The Temporal Graph Benchmark and stricter dynamic-link evaluation protocols~\cite{huang2023tgb,poursafaei2022evaluation} reinforce the importance of task-matched support and realistic temporal boundaries. A comparison based on different prediction units would not answer whether either approach is better.

In blockchain analysis, Elliptic2~\cite{bellei2024elliptic2} highlights subgraph-level supervision for anti-money-laundering. Its labels and Bitcoin-cluster units are not equivalent to the contract labels used here. Similarly, graph outlier benchmarks such as BOND and PyGOD~\cite{liu2022bond,liu2024pygod} distinguish anomaly scoring from supervised fraud classification. Historical benchmark variants are retained in supplementary diagnostics, not merged into the primary production table.

\subsection{Selective Prediction, Calibration and Risk}
SelectiveNet~\cite{geifman2019selectivenet} studies learned rejection, while early-exit systems~\cite{teerapittayanon2016branchynet} allocate computation within a model. Here escalation invokes another graph stage and still returns a prediction. This is computational deferral, not a guarantee that accepted predictions have a prescribed error rate. MC dropout~\cite{gal2016dropout} supplies an optional stochastic local predictor; the validation selection in this revision retains a single deterministic pass. We therefore make no empirical claim that MC uncertainty improves this operating point.

Post-hoc calibration~\cite{guo2017calibration} changes score semantics, and graph-specific miscalibration~\cite{hsu2022what} motivates reporting multiple calibration measures. Distribution-shift selective classification~\cite{liang2024selective} further cautions against transferring confidence rules across populations. Conformal risk control~\cite{angelopoulos2024crc}, online conformal adaptation~\cite{bhatnagar2023online}, and graph conformal prediction~\cite{zargarbashi2023conformal,zhang2025residual} provide distinct statistical frameworks. Our binomial intervals are descriptive pointwise intervals under exchangeability, not an implementation or extension of those guarantees.

\section{Data, Prediction Units and Audit Boundaries}
\subsection{Available Artifacts}
The evaluation reuses frozen training, validation and test contract graphs from the existing multi-chain reconstruction. Each bounded graph contains degree-based features, edges, a graph identifier and a contract label. The frozen models were trained previously; this cycle performs inference-interface correction and validation recalibration, not a new dataset reconstruction. Source hashes identify the retained artifacts. The original chronological transaction directory is unavailable in the current environment, so the complete upstream reconstruction and temporal preprocessing cannot be independently rerun here.

The offline unit is one contract snapshot. The replay source contains \StreamEvents{} transaction records contributed by \UniqueContracts{} contracts. Labels are inherited from contract metadata and repeated at each event. Event prevalence is \EventPrevalence{} percent. No more than thirty sampled events come from any one contract; repeated weighting is present but the workload is not dominated by a handful of high-activity contracts. The source is a chronological prefix assembled from sampled per-contract records, not the full transaction stream of any chain.

\subsection{Three Noninterchangeable Evidence Lanes}
The offline lane measures held-out contract predictions. The short-prefix lane measures warmed raw-event latency and throughput, with no fraud-positive examples on its evaluated prefix. The integrated lane measures stateful execution over the longer retrospective event sequence. They share a model interface and the primary seed's frozen policy where applicable, but not populations, graph histories or statistical interpretation. Table~\ref{tab:operating_point_identity} identifies their routing rates separately.
\input{../results/sci_v3_submission_r3/tables/table_operating_point_identity.tex}

\subsection{Temporal Qualification}
All replay events precede the latest timestamp used by the fixed training reference. In addition, \EarlyTestContracts{} held-out test contracts end before that global training boundary. Train/test membership remains disjoint at the contract level, but per-chain split membership does not prove a global temporal ordering for a pooled model. We therefore describe offline results as held-out contract estimates and replay as a retrospective systems workload. Neither constitutes a strict live temporal detection evaluation. Restoring full histories and retraining under a globally valid cutoff is required before making that stronger claim.

\section{Shared Inference and Bounded Execution}
\subsection{Local Graph and Encoder}
For each event, a chain-qualified contract key selects a bounded transaction store. The materialized graph retains at most 128 nodes and 128 edges within a ninety-day window. Each node has three features: the logarithms of one plus its incoming, outgoing and total retained degree. This representation is computed from the current retained event history; labels are not inserted into model features or live state. The local encoder is a two-layer GIN with hidden width 32, mean/max readout and optional dropout probability 0.2. For $T$ stochastic evaluations, local scores and embeddings are averaged. At $T=1$, the model runs in evaluation mode and dropout is disabled.

\subsection{Independent Relational Targets}
We compute a fixed reference of training embeddings and local scores. For each target, its eight nearest training embeddings define a bidirectional star centered on the target. Node features concatenate the embedding and raw local score. A two-layer relational GATv2 returns the target score. Targets never exchange messages with other validation, test or concurrently arriving targets. Offline batching forms disconnected stars; event-driven execution creates the identical target graph.

This interface differs from historical batched target-to-target relations and recent-event context. Frozen model weights are retained, but validation calibration is refitted on the corrected interface. This change is an execution-consistency correction, not evidence that the original relation configuration had the same predictive performance. The reference is fixed-size for a fixed training set; its memory scales with the training corpus and is not included in a claim of constant memory with respect to training-set size.

\subsection{Calibration, Routing and Final Decision}
Let $s_1,s_2$ denote raw probabilities. Separate validation-only logistic maps transform their clipped log odds:
\begin{equation}
p_j=\sigma\{a_j\operatorname{logit}(s_j)+b_j\},\qquad j\in\{1,2\}.
\end{equation}
A validation-selected local threshold $\tau_1$ defines distance $d=|p_1-\tau_1|$. A target escalates if $d\leq m$, where $m$ is a validation-derived margin quantile. On escalation,
\begin{equation}
p_f=\sigma\{w\operatorname{logit}(p_1)+(1-w)\operatorname{logit}(p_2)\}.
\end{equation}
Otherwise $p_f=p_1$. The final prediction compares $p_f$ with its separately selected threshold $\tau_f$. Thus calibration, routing and decision thresholds are not interchangeable. A direct-only control uses the calibrated local threshold; the full-relational control executes the relational stage for every target and uses a validation-selected full-fusion threshold.

The fixed validation search considers deep budgets 0.10, 0.20, 0.30 and 0.42 and fusion weights from zero to one in steps of 0.1. It maximizes validation F1, breaking ties by lower escalation and then proximity to equal weighting. Among $T\in\{1,3,5,8\}$, the smallest candidate within 0.01 validation F1 of the best is selected on the designated primary seed. Selection finishes and is persisted before test inference. The resulting $T=1$ is shared by all seeds and all event runs. Seed-specific calibrations follow the same frozen rule. The event profiler uses the primary seed's exact maps, margin, weight and thresholds; it does not fit on replay labels.

\subsection{State, Checkpointing and Memory Semantics}
The implementation limits active contract state to 5,000 chain-qualified entries, with least-recently-used eviction and per-contract temporal expiration. The embedding cache is bounded by both entries and serialized payload bytes. Synchronous work queues are drained before the next event; this tests bounded sequential execution, not sustained concurrent arrival overload. The trace writer flushes bounded chunks rather than retaining unbounded prediction dictionaries.

A checkpoint persists store, cache, active-key order, counters and random-generator states to disk. A fresh engine restores that state and reevaluates the next event; its scores and decisions are compared to uninterrupted execution. This checks durable deserialization and logical restart consistency, not fault injection under power loss or atomic external message acknowledgments. Process RSS, Python-managed allocations, CUDA allocated/reserved memory, store/cache payload sizes and trace-buffer size are sampled together. Logical payload and Python counters overlap; their sum is not asserted to partition RSS.

\section{Experimental Protocol}
\subsection{Statistical Track and Support}
Track B was recorded before final test evaluation: the five existing frozen models are used for descriptive predictive evidence, while policy-correct systems measurement takes priority over twenty-seed retraining. All comparisons use matched test contracts. Within each model, class-stratified paired bootstrap resamples provide conditional intervals for differences in F1, PR-AUC, MCC and fraud recall. Exact McNemar tests count matched prediction-error discordances and Holm adjustment controls the five-test family under the tests' assumptions. Seed differences, signs and sample standard deviations remain visible. Contracts can be statistically dependent; nominal bootstrap and binomial intervals therefore require care, and no confirmatory gain is inferred.

\subsection{Calibration and Coverage}
We report NLL, Brier score, equal-width ECE with ten and twenty bins, equal-frequency adaptive ECE, and macro one-versus-rest classwise ECE. The main table summarizes the most interpretable quantities; the supplement provides support counts and all measures. Reliability plots show actual bin counts and binomial intervals. Validation is used both for calibration and policy selection, so validation calibration quality is not an independent generalization estimate.

The descriptive risk family contains twenty-one validation-derived deep budgets. Each margin and final threshold is fixed on validation before computing test curves. Direct coverage is the fraction not escalated. Fraud-miss risk is the false-negative fraction among fraud-positive directly handled contracts. Its one-sided upper interval is exact under a binomial model, pointwise rather than simultaneous, and undefined at zero fraud-positive direct support. A dense requested grid need not yield distinct supported coverage points when scores tie or fraud support is sparse.

\subsection{Systems Measurement}
The evaluated hardware is an RTX 4070 Laptop GPU with approximately eight GiB of memory, Ubuntu 26.04.1 under WSL2, Python 3.12.13 and PyTorch 2.5.1 with CUDA 12.1. The short-prefix comparison alternates policy order over five repeats of the identical 500-event prefix after 25 warmup events. Warmup state is discarded before measurement and the random seed is reset. GPU-to-host score transfers synchronize actual model work. Instrumented wall throughput includes trace writing, memory checkpoints and restart verification; per-event latency excludes checkpoint serialization and duplicate verification work. This distinction is preserved in all tables.

Both primary and full-relational policies process the available 100,000-event source. The original histories needed for a genuine longer chronological run are missing. Repeating the prefix is not presented as 500,000 new chronological observations. No second GPU class is available, and no hardware-invariance claim is made.

\section{Results}
\subsection{Predictive and Statistical Evidence}
Table~\ref{tab:main_predictive} reports primary contract results. The observed mean F1 difference is \DeltaFOne{}, with \PositiveSeeds{} positive seed differences and \NegativeSeeds{} negative differences. Table~\ref{tab:statistical_evidence} retains per-model conditional intervals rather than replacing training variability with a pooled prediction interval. The result is descriptive under the predeclared track regardless of individual adjusted test outcomes.
\input{../results/sci_v3_submission_r3/tables/table_main_predictive.tex}
\input{../results/sci_v3_submission_r3/tables/table_statistical_evidence.tex}
\input{../results/sci_v3_submission_r3/tables/table_cross_chain_summary.tex}
Per-chain results show unequal fraud support. The Polygon test interval has no positive contracts, so its fraud discrimination scores are undefined. This table is a breakdown of a pooled trained model, not a new source-only cross-chain transfer experiment.

\subsection{Calibration and Risk--Cost Description}
\input{../results/sci_v3_submission_r3/tables/table_calibration_metrics.tex}
\begin{figure*}[t]\centering\includegraphics[width=.95\textwidth]{../results/sci_v3_submission_r3/figures/figure_reliability_r3.pdf}\caption{Primary-seed reliability with equal-frequency bins, support labels and pointwise binomial intervals. Validation and test are separated.}\end{figure*}
\begin{figure*}[t]\centering\includegraphics[width=.95\textwidth]{../results/sci_v3_submission_r3/figures/figure_risk_coverage_dense.pdf}\caption{Validation-derived descriptive test sweep. Markers are measured points; unsupported risks are omitted. Upper intervals are pointwise and do not certify a selected test operating point.}\end{figure*}
\begin{figure}[t]\centering\includegraphics[width=\columnwidth]{../results/sci_v3_submission_r3/figures/figure_accuracy_cost_frontier.pdf}\caption{Contract-lane F1 and deep-route fraction across frozen seeds. Horizontal position is an executed model-call fraction, not measured event latency.}\end{figure}
Numerical calibration results determine whether a given measure is lower; no blanket improvement claim follows from a reliability plot alone. The risk family is shown to expose the cost of leaving fraud-positive contracts on the local path. It is not used to reselect the primary policy on test labels.

\subsection{Policy-Aligned Runtime and Integrated State}
\input{../results/sci_v3_submission_r3/tables/table_primary_selective_frontier.tex}
\input{../results/sci_v3_submission_r3/tables/table_integrated_streaming.tex}
The short-prefix primary and full-relational mean latencies are \PrimaryLatency{} and \FullLatency{} milliseconds. Their difference is measured under the same warmed protocol, not inferred by multiplying routing coverage by an assumed stage cost. Short-prefix routing need not equal offline or integrated routing because the graph histories and populations differ. The supplementary identity ledger records the historical values and their incompatible configurations.

\begin{figure*}[t]\centering\includegraphics[width=.95\textwidth]{../results/sci_v3_submission_r3/figures/figure_streaming_memory_long.pdf}\caption{Retrospective replay memory. Process RSS and bounded logical store/cache payloads use separate panels. The observed horizon ends at the available source; no beyond-horizon plateau is asserted.}\end{figure*}
The integrated runs report \LossCount{} lost events and \RestartErrors{} checkpoint decision disagreements in aggregate. These counts apply to synchronous local replay. Evicted historical state is separately accounted for and is not counted as loss of the arriving event. End-to-end queue behavior under real arrival pressure and exactly-once distributed recovery remain untested.

\section{Discussion and Limitations}
\subsection{What the Evidence Supports}
The corrected interface permits a defensible within-workload comparison of selective and full relational execution. Conditional invocation is actually skipped for direct decisions, while shared calibration and fusion avoid profiling an unintended policy. Bounded store, cache and trace capacities are concrete implementation invariants. Their existence does not prove bounded total process RSS: models, references, source preloading, native libraries, serialization and allocators all contribute.

The detector's accuracy difference is less stable than a single mean suggests. Five fixed models cannot establish robust superiority, and a calibrated threshold selected on validation can change error tradeoffs on an imbalanced test set. The selected deterministic pass also means this revision does not substantiate MC uncertainty as the mechanism responsible for the measured cost difference. That optional mechanism requires a separate validation-supported evaluation before being promoted to a headline contribution.

\subsection{What Remains Unresolved}
The global training-time violation prevents a strict live temporal interpretation. Restoring the complete histories, enforcing one globally earlier training boundary and retraining all compared models is necessary to close it. A fair TGN or temporal-attention comparator additionally requires the same target contracts, labels and complete timestamped histories. The formal exclusion report explains why training it on a sampled prefix would be misleading. Its absence is a limitation, not evidence against temporal models.

No observations beyond the available replay horizon establish a memory plateau. The workload contains repeated inherited labels, not independently adjudicated event fraud. Its event-weighted F1 is therefore removed from the systems headline and retained only in diagnostic material. Even selecting each contract's last observed event produces a different snapshot from the offline benchmark. Hardware coverage, concurrency, distribution drift and labels obtained after deployment remain open evaluation dimensions.

\section{Conclusion}
DLG-StreamMC demonstrates how a bounded local-to-global graph detector can execute a validation-calibrated selective policy consistently across offline and event paths. The evidence package separates populations, exposes seed uncertainty, and measures actual policy execution and state restoration. The observed predictive change remains descriptive. Retrospective timing and configured state bounds are supported on the measured workload, while strict live temporal validity, a matched temporal baseline and longer-horizon process-memory behavior are explicitly unresolved. These boundaries are prerequisites for an honest external scientific review, not implementation details to hide in a runtime table.
\bibliographystyle{IEEEtran}
\bibliography{references}
\end{document}
